\section{Conclusion}\label{sec:conclusion}
Neural Bellman Operators connect a learned policy-specific continuation to
feasible actions and an economic payoff account. Centering identifies the
continuation errors relevant to decisions; transport describes the effect
of changing futures; full-policy verification determines whether local
accuracy extends to a dynamic economic objective. None of these operations
makes a neural representation preferable merely because its candidates
can be certified.

The new coercive residual result and its capital application complete an
all-state, full-vector-action policy certificate with continuous Gaussian
innovations and reoptimized future policies. The own-policy value and
rounding account are computed rather than replaced by a fitted loss. A
separate trainable-factor implication ties an actual neural gradient and
rank certificate to policy accuracy. The class, model-transfer, and action
implementation terms are explicit. This completes the dynamic guarantee
for the stated discrete-time economy, without asserting that the same
constants apply to nonlinear recursive utility or strategic equilibria.

All 480 frozen complete services reach the declared full-policy tolerance;
all 608 unsuccessful intermediate checks are retained. Cold and warm neural
fitting, matched conventional quadratics, and a structure-exploiting dynamic
program receive the same economic verifier. The finite randomizer makes
success probabilities and capped operation expectations well defined. The
clocks include failed checks and candidate writes. The empirical conclusions
concern those complete procedures, not the verification theorem alone.
The earlier nonlinear capital studies and every unfavorable comparison
remain part of the scientific record.

The general models of recursive utility, endogenous preferences, temporal
selves, and games continue to supply distinct economic problems for NBO.
Their original admissibility and verification accounts are retained. The
paper's method-level question remains where learning a neural continuation
improves attainable economic accuracy or total work relative to legitimate
alternatives. The completed full-policy account makes that question
sharper: an advantage must be assessed at a policy target with all its
terms supplied, rather than inferred from a successful local decision.
